% =====================================================================
% Společné pomůcky pro šablony titulní strany
% =====================================================================

% Vypnutí dělení slov -- na titulní straně působí rušivě. Kromě
% penalizace je potřeba i \emergencystretch, jinak TeX slovo rozdělí
% ve chvíli, kdy jinou přijatelnou sazbu nenajde.
\newcommand{\tpnohyphen}{%
    \hyphenpenalty=10000\exhyphenpenalty=10000%
    \tolerance=9999\emergencystretch=3em\relax
}

% Školní rok. Čísla se sázejí v matematickém režimu, takže barvu je
% nutné nastavit i pro "math text" -- jinak zůstanou tmavě modrá.
% #1 = barva
\newcommand{\tpdate}[1]{{\setbeamercolor{math text}{fg=#1}\color{#1}\insertdate}}

% Nastavení textu vysázeného na barevném (tmavém) podkladu:
% matematika i makro \Astar musí být bílé, ne černé/tmavě modré.
\makeatletter
\newcommand{\tpinverse}{%
    \tpnohyphen
    \setbeamercolor{math text}{fg=white}%
    \@ifundefined{Astar}{}{\renewcommand{\Astar}{\textup{A}\textsuperscript{*}}}%
}
\makeatother

% Téma přednášky (podtitul) -- jednotné sázení ve všech variantách.
% Pozor: \tpnohyphen musí být mimo skupinu s fontem -- parametry
% odstavce se čtou až při jeho ukončení. Nastavení \hyphenchar na -1
% dělení vypne i tam, kde by TeX jinak neměl jiné řešení.
\newcommand{\tptopic}[1]{%
    \tpnohyphen
    {\usebeamerfont{title}\Large
     \hyphenchar\font=-1\relax
     #1\par
     \hyphenchar\font=`\-\relax}%
}
